\begin{folhadeaprovacao}

  \begin{center}
    {\ABNTEXchapterfont\large\imprimirautor}

    \vspace*{\fill}\vspace*{\fill}
    {\ABNTEXchapterfont\bfseries\Large\imprimirtitulo}
    \vspace*{\fill}

    \hspace{.45\textwidth}
    \begin{minipage}{.5\textwidth}
        \imprimirpreambulo
    \end{minipage}%
    \vspace*{\fill}
   \end{center}

   Trabalho aprovado. \imprimirlocal, \today:

   \assinatura{\textbf{\imprimirorientador} \\ Orientador}
   \assinatura{\textbf{Membro da banca} \\ A definir}
   \assinatura{\textbf{Membro da banca} \\ A definir}

   \begin{center}
    \vspace*{0.5cm}
    {\large\imprimirlocal}
    \par
    {\large\imprimirdata}
    \vspace*{1cm}
  \end{center}

\end{folhadeaprovacao}

\begin{agradecimentos}
Este manuscrito registra a fundamentação inicial da pesquisa de mestrado junto ao PPGCC/UFSJ. Agradeço ao orientador Michel Carlo Rodrigues Leles pelo acompanhamento e à Universidade Federal de São João del-Rei pelas condições de trabalho.
\end{agradecimentos}

\begin{resumo}

Notícias financeiras em português carregam informação que o mercado pode precificar, mas transformar esse texto em uma série temporal reprodutível ainda é um problema aberto --- sobretudo em setores pouco cobertos pela literatura de processamento de linguagem natural. Esta dissertação fundamenta a proposta de um Índice Temporal Informacional (ITI) a partir de notícias corporativas de empresas brasileiras de saneamento de capital aberto, descreve o protocolo operacional já montado no laboratório e registra andamentos experimentais sem fechá-los como resultado. O ITI não replica um índice publicado: toma a cadeia metodológica ``texto agregado no tempo e validado externamente'', o escore de sentimento de domínio em português e a tradição de comparar medidas simples antes de construções mais elaboradas. A lacuna é a ausência de um microíndice empresa a empresa, baseado em notícias, para o saneamento --- setor marcado pelo Marco Legal e por eventos institucionais como a privatização da Sabesp. A avaliação de desempenho permanece em aberto, condicionada à qualidade do classificador no domínio.

\vspace{\onelineskip}

\noindent
\textbf{Palavras-chave}: sentimento de notícias; índice temporal; saneamento; processamento de linguagem natural; mercado brasileiro.
\end{resumo}

\begin{resumo}[Abstract]
\begin{otherlanguage*}{english}

Financial news in Portuguese can carry information that markets may price, yet turning that text into a reproducible time series remains an open problem --- especially in sectors that NLP research has barely covered. This dissertation motivates an Informational Temporal Index (ITI) built from corporate news about publicly listed Brazilian sanitation firms, describes the operational protocol already in place in the laboratory, and records preliminary experiments without treating them as final results. The ITI is not a copy of a published index: it borrows the methodological chain of aggregating text over time and validating it externally, a Portuguese domain sentiment score, and the habit of comparing simple baselines before a more elaborate construction. The gap is the lack of a firm-level, news-based index for sanitation --- a sector shaped by the Sanitation Legal Framework and institutional events such as Sabesp's privatization. Performance evaluation remains open, conditional on classifier quality in this domain.

\vspace{\onelineskip}

\noindent
\textbf{Key-words}: news sentiment; temporal index; sanitation; natural language processing; Brazilian market.
\end{otherlanguage*}
\end{resumo}
